\section{Decoding with a hidden coordinate}
\label{sec:decoding}

We now convert the Fourier mass supplied by Proposition~\ref{prop:extraction}
into a labeling of the original affine game.  The difficulty is that the
mass is measured after fixing a row context.  A labeler at an $X$-vertex
does not know the projection on its incident edge, and therefore cannot
usually lift that context to its own label space.  We resolve this by
repeating the game coordinates and planting one coordinate on which the
context and the latent resampling rows have zero linear coefficients.
Hiding the planted coordinate makes the resulting
distribution close to the original one.
The comparison across the unknown projection then involves only the $2s$
free shifts, so its loss is independent of $n$ and $T$.

Throughout this section, the cut tables are fixed.  We use the tree
parameters and the row coordinates of the preceding sections.  For each
choice of a slice $\mathcal D$ and a node $v\in\mathcal D$, fix the
coordinate change
\[
 R\longleftrightarrow (L,\Gamma),\qquad
 L=(L_{gb})_{g\in[s],\,b\in\{0,1\}}\in H^{2s},
 \quad \Gamma\in H^{k-2s},
\]
by choosing its representative leaves in a deterministic order.  This
change uses only additions of row coordinates and depends only on the
tree and $v$.  In particular, it commutes with pullback of affine forms.
The masses $m_q$ always refer to these coordinates and to the uniform
average over the slice tables in Equation~\eqref{eq:mass}.

\subsection{Removing the row smoothing}

Let $R^{\mathrm{new}}_1,\ldots,R^{\mathrm{new}}_k$ be independent uniform
elements of $H$, and independently let $B_1,\ldots,B_k$ be Bernoulli
variables with $\Pr(B_i=1)=1-r$.  Set
\[
 E_i=B_iR^{\mathrm{new}}_i,
 \qquad E=(E_1,\ldots,E_k).
\]
For any fixed row array $R$, the array $R+E$ has the distribution of
the one-copy row resampling used by the test.  Write $(E_L,E_\Gamma)$
for the coordinates of $E$ under the above linear change of variables.

\begin{lemma}[Removing smoothing]\label{lem:unsmooth}
For every fixed $\underline y$, $\mathcal D$, $v$, $\Gamma$,
$F_{\mathrm{out}}$, and label $q\in\mathcal A$,
\begin{equation}\label{eq:unsmooth}
 m_q(\bar G_{\underline y};\Gamma,F_{\mathrm{out}})
 \leq
 \E_{E,\,\underline e\mid\underline y}
 m_q(G_{\underline e};\Gamma+E_\Gamma,F_{\mathrm{out}}).
\end{equation}
Here the noise $E$ is independent of the conditionally sampled edge
tuple $\underline e$.
\end{lemma}

\begin{proof}
For the fixed context, regard a function of $(L,F_{\mathcal D})$ as an
element of the real Hilbert space with the uniform probability measure.
Let $\Pi_q$ be the orthogonal projection onto the shift Fourier indices
$Q$ satisfying $Q_{g0}+Q_{g1}=q$ for every $g$.  Then the corresponding
mass is the squared norm of this projection.  The identity
\[
 \bar G_{\underline y}(L,\Gamma,F)
 =\E_{E,\,\underline e\mid\underline y}
 G_{\underline e}(L+E_L,\Gamma+E_\Gamma,F)
\]
and convexity of the squared Hilbert-space norm give the desired
inequality.  Indeed, translation by $E_L$ commutes with $\Pi_q$ and
preserves its norm: on each shift Fourier coefficient it only
multiplies by a sign.
\end{proof}

Choose once and for all an ordering of each label space, and define
\begin{equation}\label{eq:decoder-choice}
 q^*=q^*(\underline y,\mathcal D,v,\Gamma,F_{\mathrm{out}})
 \in\operatorname*{arg\,max}_{q\in\mathcal A}
 m_q(\bar G_{\underline y};\Gamma,F_{\mathrm{out}})
\end{equation}
using that ordering to break ties.  This definition uses the original
conditional edge average and the original row-smoothing operator.
It will remain unchanged when we modify the distribution of the row
data below.

Sample $\underline y$, the slice, its uniform node, the context, and
the outside tables as in Proposition~\ref{prop:extraction}, and then
sample $R^{\mathrm{new}}$, $B$, and $\underline e\mid\underline y$.
Combining that Proposition with Lemma~\ref{lem:unsmooth} gives
\begin{equation}\label{eq:unplanted-score}
 \E\,m_{q^*}(G_{\underline e};
              \Gamma+E_\Gamma,F_{\mathrm{out}})
 \geq p_*\theta.
\end{equation}
The mass in this expectation lies in $[0,1]$, since $G_{\underline e}$
is bounded in absolute value by one.

\subsection{Planting a coordinate}

Write each affine row in coordinates on
$\mathcal A=(\F^n)^T$.  Apart from its single constant coefficient, it
has one block of $n$ linear coefficients for each of the $T$ slots.
For a slot $a\in[T]$, let $Z_a$ be the event that the slot-$a$ blocks
vanish in all $k-2s$ context rows and all $k$ latent rows
$R^{\mathrm{new}}_i$.  No condition is imposed on any constant
coefficient.  Under the original, unplanted law,
\begin{equation}\label{eq:zero-slot-probability}
 p_0:=\Pr(Z_a)=2^{-n(2k-2s)}.
\end{equation}
Conditional on the depth, slice, and node, the events $Z_a$ are
independent: they concern disjoint blocks of independent uniform
coefficients.

The \emph{planted experiment} chooses a uniform mark $a_*\in[T]$ and
sets these blocks to zero at $a_*$.  All other row coefficients, the
noise masks, the outside tables, and the game edges retain their
original distributions.  Equivalently, conditional on $a_*=a$, this
is the unplanted law conditioned on $Z_a$.

\begin{lemma}[Hiding the planted coordinate]\label{lem:planting}
Let $\mathsf U$ be the unplanted law and $\mathsf P$ the planted law
after forgetting the mark.  Both laws are taken on the full data
\[
 \bigl(\underline e,J_0,\mathcal D,v,F_{\mathrm{out}},
        \Gamma,R^{\mathrm{new}},B\bigr).
\]
Then
\begin{equation}\label{eq:planting-tv}
 \TV(\mathsf P,\mathsf U)
 \leq \frac12\sqrt{\frac{1-p_0}{Tp_0}}
 \leq \frac{1}{2\sqrt{Tp_0}}.
\end{equation}
Consequently, if
\begin{equation}\label{eq:decoder-repetition}
 T\geq\frac{1}{p_0(p_*\theta)^2},
\end{equation}
then under the planted law,
\begin{equation}\label{eq:planted-score}
 \E_{\mathsf P}\,m_{q^*}(G_{\underline e};
               \Gamma+E_\Gamma,F_{\mathrm{out}})
 \geq \frac{p_*\theta}{2}.
\end{equation}
\end{lemma}

\begin{proof}
Let $Z=\sum_{a=1}^T\one_{Z_a}$.  Averaging the conditional densities
over the uniform mark gives
\[
 \frac{d\mathsf P}{d\mathsf U}=\frac{Z}{Tp_0}.
\]
Conditional on the tree data, $Z$ has the binomial distribution with
parameters $T,p_0$; these parameters do not depend on the tree data.
Thus Cauchy--Schwarz gives
\[
 \TV(\mathsf P,\mathsf U)
 =\frac12\E_{\mathsf U}\abs{\frac{Z}{Tp_0}-1}
 \leq\frac12\sqrt{\frac{1-p_0}{Tp_0}}.
\]
The score in Equation~\eqref{eq:unplanted-score} is a $[0,1]$-valued
function of the displayed data without the mark.  In particular,
$q^*$ is the fixed function defined in Equation~\eqref{eq:decoder-choice};
it does not use the mark, the sampled incident edges beyond their
$Y$-endpoints, or the sampled noise.  Its statistical dependence on
zero slots causes no difficulty.  Expectations of this score under
the two laws differ by at most their total variation distance.
Equations~\eqref{eq:unplanted-score} and
\eqref{eq:decoder-repetition} therefore imply
Equation~\eqref{eq:planted-score}.
\end{proof}

\subsection{Fourier mass under an affine restriction}

The next elementary observation accounts for the loss in decoding.
Let $\rho:\widetilde{\mathcal A}\to\mathcal A$ be a surjective affine
map with two-element fibers, and write
\[
 \widetilde H=\Aff(\widetilde{\mathcal A},\F),\qquad
 \iota=\rho^*:H\hookrightarrow\widetilde H.
\]
Pullback is linear on the vector spaces of affine forms, even when
$\rho$ has a nonzero translation part.  The image of $\iota$ has
codimension one, and $\iota(\one)=\one$.
Its dual restriction map
$\iota^*:\widetilde H^*\to H^*$ has two-element fibers and, on
evaluation functionals, agrees with $\rho$.

\begin{lemma}[Restriction of shift mass]\label{lem:restriction}
Fix a context and the outside tables.  Let
$\widetilde f(\widetilde L,F_{\mathcal D})$ be a function on
$\widetilde H^{2s}$ and the slice tables.  Suppose its shift Fourier
support has the common-pair-sum property
\[
 \widetilde Q_{g0}+\widetilde Q_{g1}
 =\widetilde Q_{h0}+\widetilde Q_{h1}
 \quad\text{for all }g,h\in[s].
\]
Set $f(L,F_{\mathcal D})=
\widetilde f(\iota L,F_{\mathcal D})$, where $\iota$ acts on each
shift.  For every label $q\in\mathcal A$,
\begin{equation}\label{eq:restriction-mass}
 m_q(f)\leq 2^{2s}
       \sum_{\substack{\widetilde q\in\widetilde{\mathcal A}\\
                       \rho(\widetilde q)=q}}
          \widetilde m_{\widetilde q}(\widetilde f).
\end{equation}
The masses on both sides use the common-pair-sum definition of
Equation~\eqref{eq:mass}, with the respective affine row spaces.
\end{lemma}

\begin{proof}
For every fixed choice of slice tables, Fourier inversion and
character orthogonality give
\[
 \widehat f(Q)=
 \sum_{\substack{\widetilde Q\\
                  \iota^*\widetilde Q_{gb}=Q_{gb}\ \forall g,b}}
       \widehat{\widetilde f}(\widetilde Q).
\]
The sum has $2^{2s}$ indices.  Hence its squared absolute value is at
most $2^{2s}$ times the sum of the squared absolute values of its
terms.  Now sum over all $Q$ with pair sums $q$.  Every nonzero
coefficient on the right has a common pair sum $\widetilde q$ whose
restriction is $q$.  Moreover,
\[
 \widetilde q(\one)=\widetilde q(\iota(\one))=q(\one)=1,
\]
so $\widetilde q$ is an evaluation functional, hence a label in
$\widetilde{\mathcal A}$, and $\rho(\widetilde q)=q$.
Each lifted index is counted only once.  Averaging over the slice
tables proves Equation~\eqref{eq:restriction-mass}.
\end{proof}

\subsection{The endpoint labelers}

Equation~\eqref{eq:planted-score} is an average over full edge tuples.
To turn it into a labeling of the original game, we place the input edge
at the planted coordinate and make all other edges public.  The zero
coefficients ensure that the needed context lift does not require the
projection on the input edge.

\begin{proposition}[Decoding]\label{prop:decoding}
Suppose the extraction conclusion of Proposition~\ref{prop:extraction}
holds with constants $p_*,\theta>0$.  If $T$ satisfies
Equation~\eqref{eq:decoder-repetition}, then the original affine
game has value at least
\begin{equation}\label{eq:decoded-value}
 \Val(\mathcal G)\geq 2^{-2s}\frac{p_*\theta}{2}.
\end{equation}
\end{proposition}

\begin{proof}
We construct randomized labelers that use only their own endpoint
and shared randomness.  The entire game and the fixed cut tables
are available as background data.  No efficiency claim about the
labelers is needed to lower-bound game value.

\paragraph{Shared data and the $Y$-labeler.}
Given a random input edge $e=(x,y)$ from the original edge law,
choose a uniform public mark $a_*\in[T]$.  At every other slot,
publicly sample an independent edge, including its endpoints and
projection.  Insert the input edge at the marked slot.  Its
projection is used only in the analysis; it is not revealed to
either labeler.  Conditional on any mark, the resulting full edge
tuple has the product edge law; hence the tuple is independent of
the mark.  Each labeler knows its own endpoint tuple, and both know
all edges outside the marked slot.

Publicly sample the depth, slice, node, outside tables, context,
latent rows, and masks according to the planted experiment above.
The row sampling uses the fixed alphabet dimensions, the public
tree coordinates, and the mark, and is independent of the input edge.
Together with the edge tuple, these data have exactly the planted
law used in Equation~\eqref{eq:planted-score}.  The $Y$-labeler
forms its endpoint tuple $\underline y$, computes the label $q^*$
from Equation~\eqref{eq:decoder-choice}, and outputs its marked
component $q^*_{a_*}$.  All quantities used in that computation are
available to this labeler: in particular the conditional average
defining $\bar G_{\underline y}$ can be computed from the game and
cut as fixed background data.  It does not use the actual marked
edge or its projection.

\paragraph{The $X$-labeler.}
To define the $X$-labeler, replace the marked factor of $\mathcal A$
by the $X$-alphabet:
\[
 \widetilde{\mathcal A}
 =\prod_{a\ne a_*}\F^n\ \times\ \F^{n+1},
\]
with factors in their original slot order.  Let
$\mathcal B=(\F^{n+1})^T$.  The full projection factors as
\[
 \mathcal B\xrightarrow{\ \kappa\ }
 \widetilde{\mathcal A}\xrightarrow{\ \rho\ }\mathcal A.
\]
Here $\kappa$ is the public edge projection at every unmarked slot
and the identity at the marked slot; $\rho$ is the identity at
every unmarked slot and the unknown map $\pi_e$ at the marked
slot.  Thus $\kappa$ is known to the $X$-labeler, while $\rho$ is
used only in the analysis.

The planted zeros imply that every row in $\Gamma$ and
$R^{\mathrm{new}}$ has zero linear coefficient block in the marked slot.
The same holds for $E$, $E_\Gamma$, and
$\Gamma'=\Gamma+E_\Gamma$, since the noise masks and context change
only multiply by scalars and add rows.  Consequently
$\widetilde\Gamma'=\rho^*\Gamma'$ is known to the $X$-labeler:
it is obtained by leaving every unmarked coefficient and constant
coefficient unchanged and setting all marked coefficients to zero.
This description does not require $\pi_e$, including its affine
translation part.

The $X$-labeler forms the bounded raw function
\[
 \widetilde G(\widetilde R,F)
 =\sigma^o_{\underline x}
   \bigl(W_{\widetilde{\mathcal A}}(\widetilde R,F)
                \circ\kappa\bigr),
 \qquad \widetilde R\in\widetilde H^k.
\]
Fix its context to $\widetilde\Gamma'$ and its outside tables to
the public $F_{\mathrm{out}}$, leaving all $2s$ shifts free in
$\widetilde H$.  Let
$\widetilde m_{\widetilde q}$ be its masses.  They are nonnegative
and their sum over labels is at most one, by Parseval and
$\abs{\widetilde G}\leq1$.  The labeler samples a label
$\widetilde q$ with probability at least
$\widetilde m_{\widetilde q}$ for each label, assigning any leftover
probability to a fixed label.  It uses an additional independent
public uniform random number to sample from this distribution and
outputs the marked component of $\widetilde q$.  This sampling does
not use $q^*$.

\paragraph{Comparison on the sampled edge.}
We verify the restriction hypothesis needed to compare these
probabilities with the score.  For any $h\in\widetilde H$, shifting
both child generators in one pair by $h$ multiplies the output
of gate $v$ by $(-1)^{h(\widetilde l)}$ at every label
$\widetilde l$.  Applying this shift in two distinct pairs leaves
that output, and hence the whole word, unchanged.  These operations
only translate the free shifts and preserve the context.
Character orthogonality therefore shows that every shift Fourier
index of $\widetilde G$ has a common pair sum.  This argument holds
for each fixed choice of all gate tables.

Word generation commutes with pullback of rows, pointwise in the
label.  Since $\pi_{\underline e}=\rho\circ\kappa$, restricting
each free shift of $\widetilde G$ to $\rho^*H$ yields exactly
$G_{\underline e}$ with context $\Gamma'$ and the same tables.
The map $\rho$ has two-element fibers, so Lemma~\ref{lem:restriction}
applies and gives
\[
 \sum_{\rho(\widetilde q)=q^*}\widetilde m_{\widetilde q}
 \geq 2^{-2s}
 m_{q^*}(G_{\underline e};\Gamma',F_{\mathrm{out}}).
\]
Conditional on all the planted data, the left side lower-bounds
the probability that the sampled full label projects to $q^*$.
That event in particular implies satisfaction of the marked
input edge.  Averaging and using Equation~\eqref{eq:planted-score}
gives success probability at least $2^{-2s}p_*\theta/2$.

Finally, fix all shared randomness, including the uniform number
used to sample the $X$-label.  The same fixed data define the two
strategies at every endpoint, so they give a deterministic labeling
of the original game.  Their average value over the shared
randomness has the lower bound just proved, so some fixed choice
attains it.  This proves Equation~\eqref{eq:decoded-value}.
\end{proof}

The decoding loss depends on $s,p_*,\theta$, and hence on the target
cut gap, but not on the alphabet dimension or the repetition
length.  After choosing a game soundness smaller than the right
side of Equation~\eqref{eq:decoded-value}, one may therefore fix
the alphabet dimension and then choose $T$ by
Equation~\eqref{eq:decoder-repetition}.
